% Ce chapitre repose sur :
%   DOC: (aucun)
%   OBS: apport-personnel
%   HYP: (aucun)
%   CHI: conclusions-avec-relation, conclusions-sans-relation, exigences-calculables, exigences-champs-manquants, exigences-non-calculables, exigences-rubriques
%   CHI: exigences-sans-objet, fichiers-de-controle, limites-des-faits, modeles-total, relations-injectees, relations-non-reprises
%   CHI: releve-champs-non-employes, releve-champs-total, taches-graphe, tdb-indicateurs, tdb-pages
%   SER: (aucune)
%   SECTIONS: 0

\chapter*{Conclusion générale}
\addcontentsline{toc}{chapter}{Conclusion générale}
% Un chapitre étoilé ne met pas à jour la marque de l'en-tête courant : sans cette ligne,
% les pages de ce texte portent l'en-tête de ce qui le précède — mesuré, « TABLE DES
% MATIÈRES » sur la deuxième page de l'introduction et « CHAPITRE 9 » sur celle de la
% conclusion. Le défaut ne se voit qu'en regardant le document composé.
\markboth{Conclusion générale}{Conclusion générale}

\section*{Ce qui a été fait}

Une chaîne de données complète a été construite, du fichier déposé jusqu'à l'écran :
\chiffre{modeles-total} modèles de transformation, \chiffre{taches-graphe} tâches orchestrées, un
rapprochement probabiliste d'identités évalué contre une vérité terrain connue, et un tableau de
bord de \chiffre{tdb-pages} pages portant \chiffre{tdb-indicateurs} indicateurs. Elle se rejoue en
entier, et \chiffre{fichiers-de-controle} fichiers de contrôle la surveillent.

Ce n'est pas le volume de ce travail qui en fait l'intérêt, mais \textbf{ce qu'il refuse
d'affirmer}. Chaque nombre du rapport vient d'une commande qui le produit ; chaque affirmation porte
l'étiquette de ce qui l'étaye — un document, une observation, ou une convention posée ; chaque
relation injectée est déclarée, et le rapport dit laquelle des conclusions la reproduit.
\textbf{Le résultat le plus directement utilisable est la correspondance avec ce que la
nomenclature nationale attend} : sur \chiffre{exigences-rubriques} rubriques attendues d'un
établissement nommé, \chiffre{exigences-calculables} sont calculables par la chaîne.

\section*{Ce qui n'a pas été fait, et pourquoi}

\textbf{Des rubriques de la nomenclature restent hors d'atteinte, et aucune ne l'est par manque de
temps.} Elles sont \chiffre{exigences-non-calculables}, bloquées chacune par un champ absent du
système d'information, et \chiffre{exigences-champs-manquants} champs suffiraient à les débloquer.

Il s'y ajoute \chiffre{exigences-sans-objet} rubriques sans objet, et la distinction compte : elles
ne manquent pas au logiciel, c'est l'établissement qui n'exerce pas l'activité.

\textbf{Le modèle dimensionnel porte \chiffre{limites-des-faits} limites documentées plutôt que
contournées}, et toutes tiennent à la même cause : la source ne porte pas ce qu'il faudrait pour
aller plus loin. Fabriquer la distinction manquante aurait produit une donnée qu'aucune ligne
n'établit.

\textbf{Sur \chiffre{relations-injectees} relations injectées,
\chiffre{relations-non-reprises} ne sont reprises par aucune conclusion}, le plus souvent parce que
la colonne qui les porterait n'est calculée dans aucune couche aval. Le régime de couverture, l'âge
et l'adressage existent dans la source et s'arrêtent en chemin.

\textbf{Le relevé d'observation, lui, est employé en entier ou presque.} Sur
\chiffre{releve-champs-total} éléments relevés au poste,
\chiffre{releve-champs-non-employes} seulement ne sont cités nulle part — et ce sont tous des
boutons. Aucun champ de donnée observé n'est resté inemployé.

\textbf{Enfin, aucune conclusion du chapitre d'analyse n'établit un fait sur le service.} Sur les
vingt-deux qu'il porte, \chiffre{conclusions-avec-relation} reproduisent une relation injectée et
\chiffre{conclusions-sans-relation} ne viennent d'aucune. Ce n'est pas un défaut d'exécution :
c'est la conséquence nécessaire d'un jeu produit.

\section*{Les pistes qui n'ont pas été retenues}

Trois pistes ont été écartées, et les trois pour le même genre de raison. \textbf{Inventer les
libellés manquants} : un code nu est illisible et le montre, un libellé inventé est lisible et cache
qu'il est inventé. \textbf{Régénérer le jeu de données} pour effacer le désaccord des deux faits sur
la fin d'un épisode : un jeu qu'on régénère jusqu'à ce qu'il soit propre n'apprend plus rien.
\textbf{Illustrer le tableau de bord par des captures d'écran} : le dépôt interdit toute image du
système d'information.

\section*{Ce qu'il faudrait pour passer à des données réelles}

La chaîne est écrite pour lire des fichiers, non pour lire ce générateur : la question n'est pas
de savoir si elle fonctionnerait, mais \textbf{ce qui changerait de nature}. Trois choses.

\textbf{La vérité terrain disparaît.} Précision et rappel se calculent aujourd'hui parce qu'on sait
lesquels des doublons sont vrais ; sur des données réelles, cette connaissance n'existe pas.
L'évaluation changerait donc de nature — il faudrait un échantillon annoté à la main, avec son coût
et son incertitude propres — et les grandeurs publiées au chapitre du rapprochement cesseraient
d'être mesurables telles quelles.

\textbf{Les distributions ne sont plus posées.} Délais, taux d'absentéisme, pyramide de gravité aux
urgences sont aujourd'hui calibrés sur des sources publiées puis injectés. Sur des données réelles,
les contrôles qui vérifient que la chaîne retrouve le paramètre injecté n'auraient plus de second
membre, et devraient céder la place à des contrôles de forme, plus faibles et à écrire.

\textbf{Les relations ne sont plus injectées, et c'est le changement décisif : un résultat cesse de
démontrer et se met à établir.} Aujourd'hui, retrouver que l'absentéisme croît avec le délai à
l'intérieur d'une spécialité ne prouve rien d'autre que la capacité de la chaîne à retrouver ce
qu'on y a mis. Sur des données réelles, le même calcul deviendrait un constat sur le service — avec
tout ce que cela impose de prudence supplémentaire sur la population, le dénominateur et la période.

Trois conditions matérielles s'y ajoutent. Une convention d'accès aux données, que la contrainte de
confidentialité rend nécessaire et qui n'est pas du ressort d'un stage. Une extraction périodique
depuis le système d'origine, dont le format et la cadence restent à établir. Et les champs manquants nommés plus haut, qui
cesseraient d'être une limite documentée pour devenir un obstacle à lever.

\textbf{Apport personnel.} Le paragraphe qui suit relève de la rédaction personnelle et ne peut
être écrit que par l'auteur du travail.

\aRediger{apport-personnel}{ce que ce travail a apporté personnellement : les compétences acquises,
les difficultés rencontrées et la façon dont elles ont été surmontées, et ce que l'expérience du
service a changé dans la manière d'aborder un problème de
données}\releve{apport-personnel}
